\section{Introducción: ¿por qué prestar atención a MoE?}

Esta es la primera clase de la serie «Arquitectura de modelos de lenguaje grandes». En los últimos años, la evolución de las arquitecturas de los modelos grandes ha sido vertiginosa, y entre ellas la arquitectura \textbf{Mixture of Experts (MoE)} ha atraído mucha atención por su enorme ventaja en eficiencia de inferencia.

Tomemos como ejemplo DeepSeek R1: es un modelo MoE con alrededor de 670,000 millones de parámetros totales, pero en la inferencia solo activa alrededor de 72,000 millones de parámetros. En comparación con un modelo Dense completo de 72,000 millones de parámetros, la velocidad de inferencia de ambos es básicamente la misma; esto significa que un modelo MoE, al mantener una capacidad de modelo enorme (gran capacidad), logra una eficiencia de inferencia muy alta (pocos parámetros activados).

\begin{importantbox}{El valor central de MoE}
\begin{itemize}
  \item \textbf{Un número total de parámetros grande} $\Rightarrow$ gran capacidad del modelo, gran capacidad de desempeño
  \item \textbf{Un número de parámetros activados pequeño} $\Rightarrow$ inferencia rápida, alta eficiencia
  \item En la época de Qwen2.5 todavía se usaban modelos totalmente Dense; para Qwen3 ya se adoptan ampliamente variantes MoE
\end{itemize}
\end{importantbox}

En esta clase se toman \textbf{Qwen3-32B} (Dense) y \textbf{Qwen3-30B-A3B} (MoE) como objeto de estudio, y a través de comparar y calcular su número de parámetros se busca entender la arquitectura MoE.
